% This suppresses warnings about most overfull hboxes.
% Clearly it would be better to fix them, but if we anyway want to
% change the layout that can wait
\hfuzz=50pt
\usepackage{fancyhdr}

% The general view was that parskip without indent was most readable, so use it
\usepackage{parskip}

% Uses some packages:
\usepackage{amsmath}
\allowdisplaybreaks

% For non-floating figures:
\usepackage{float}


% Prefer pdf for pdf
% https://tex.stackexchange.com/questions/1072/which-graphics-formats-can-be-included-in-documents-processed-by-latex-or-pdflat
\usepackage{ifpdf}
\ifpdf
\makeatletter
\let\orig@Gin@extensions\Gin@extensions
\def\Gin@extensions{.pdf,\orig@Gin@extensions} %prepend .pdf before .png
\makeatother
\fi

% Longer section numbers in toc
\ifpdf
\usepackage{tocloft}
\setlength{\cftsecnumwidth}{3em}
\fi

% Need more text:
\usepackage[margin=1in]{geometry}

%\usepackage{t1enc}
\usepackage{graphicx}
\graphicspath{{media/}{../media/}} % so that chapter files can also find the images
\def\figureautorefname{figure} % Use lower case initial, just like for sections.

% Don't use subfig.sty until this LaTeXML issue has been fixed:
% - https://github.com/brucemiller/LaTeXML/issues/1292
% Leaving commented-out code as a hint to anyone who considers using subfig.
%\usepackage{subfig}
%\newcommand{\subfigureautorefname}{\figureautorefname}

\usepackage{verbatim}
% The fixltx2e that was used for textsubscript is no longer needed for pdf, but is needed by latexml
\ifpdf
\else
\usepackage{fixltx2e} %for textsubscript
\fi

% Allow lot of figures on page:
\renewcommand{\textfraction}{0}
\renewcommand{\topfraction}{1}
\renewcommand{\bottomfraction}{1}

% Some problem with raggedright for latexml
\ifpdf
\renewcommand{\newline}{\hspace*{\fill}\linebreak}
\else
\renewcommand{\newline}{\hspace*{\fill}\raggedright\linebreak}
\fi
\setcounter{totalnumber}{1000}
\raggedbottom
\usepackage{hyperref}
% The bookmarks are used as table-of-contents in Acrobat
% The default (magenta for urls and red for internal links) isn't nice
\hypersetup{bookmarksnumbered=true, colorlinks=true, urlcolor=blue, linkcolor=blue}
% Using longtable to have tables split across pages
% Most starting with \hline \endhead to avoid underfull vbox and not have first part duplicated on new pages
%
% But in some cases there are actual headings that are copied to next page.
% Note that this can often lead to overfull vbox
% See https://tex.stackexchange.com/questions/71096/underfull-vbox-in-each-longtable-can-it-be-fixed-or-must-it-be-ignored
\usepackage{longtable}
\usepackage{multirow} % for multirow entries in tables
\usepackage{listings}
%\usepackage{enumitem} % Package not available for continuous integration builds.
\usepackage{color}
\usepackage[table]{xcolor}
%\def\doublelabel#1{\label{#1}\hypertarget{#1}{}}
\def\doublelabel#1{\label{#1}}
\def\S0#1#2{\hypertarget{#2}\chapter{#1}\label{#2}}
\def\S1#1#2{\hypertarget{#2}\section{#1}\label{#2}}
\def\S2#1#2{\hypertarget{#2}\subsection{#1}\label{#2}}
\def\S3#1#2{\hypertarget{#2}\subsubsection{#1}\label{#2}}
\def\S4#1#2{\hypertarget{#2}\paragraph{#1}\label{#2}}
\def\S5#1#2{\hypertarget{#2}\subparagraph{#1}\label{#2}}

% Simple numbering of lists, should preferably use \usepackage{enumitem}
% https://stackoverflow.com/questions/2007627/latex-how-can-i-create-nested-lists-which-look-this-1-1-1-1-1-1-1-2-1-2
\renewcommand{\labelenumi}{\arabic{enumi}.}
\renewcommand{\labelenumii}{\labelenumi\arabic{enumii}.}
\renewcommand{\labelenumiii}{\labelenumii\arabic{enumiii}.}
\renewcommand{\labelenumiv}{\labelenumiii\arabic{enumiv}.}

% Non-normative content without paragraph breaks at beginning and end.
% In most cases, one should use the 'nonnormative' environment instead.
\newenvironment{nonnormative*}[0]{%
{[}\itshape\ignorespaces
}{%
\unskip\upshape{]}%
}

% Non-normative content with paragraph breaks at beginning and end.
\newenvironment{nonnormative}[0]{%
\par\begin{nonnormative*}%
}{%
\end{nonnormative*}\par%
}

% Example environment, as a special case of non-normative content.
\newenvironment{example*}[1][\unskip]{%
\begin{nonnormative*}Example #1:
}{%
\end{nonnormative*}
}
\newenvironment{example}[1][\unskip]{%
\begin{nonnormative}Example #1:
}{%
\end{nonnormative}
}

% Like \emph, but marking this as the place where new terminology is introduced.
% In the future, we can change this to something that will automatically add this
% to an index/glossary.
\newcommand{\firstuse}[1]{\emph{#1}}

% Formatting for filenames, URIs, etc.
% For now, the content is wrapped in \mbox to avoid hyphenation to inject
% hyphens.  A better solution could be to change the directory separator
% characters to also produce an \allowbreak{} following the directory separator.
\newcommand{\filename}[1]{\mbox{\textsf{#1}}}

% Formatting of table headings.
\newcommand{\tablehead}[1]{\textit{#1}}

\newcommand{\autonameref}[1]{\autoref{#1}~\emph{\nameref*{#1}}}

\newcommand{\glossaryitem}[1]{\textbf{#1}}

\newcommand{\bibitemtitle}[1]{\emph{#1}}

% Environment for definitions.
\usepackage{amsthm}
\newtheoremstyle{mlsdefinition}
  {\topsep}   % ABOVESPACE
  {\topsep}   % BELOWSPACE
  {}          % BODYFONT
  {}          % INDENT (empty value is the same as 0pt)
  {\bfseries} % HEADFONT
  {.}         % HEADPUNCT
  {.75em}     % HEADSPACE
  {Definition~#2. \firstuse{#3}}        % CUSTOM-HEAD-SPEC
\theoremstyle{mlsdefinition}
% In order to show where the definition ends, we put a \qed at the end.  This can't be done using \newtheoremstyle,
% so instead we make the 'definition' environment a wrapper around the amsthm-base environment.
\makeatletter
\newtheorem{definition@inner}{Definition}[chapter]
\newcommand{\definition@innerautorefname}{definition}
\newenvironment{definition}{\begin{definition@inner}}{\qed\end{definition@inner}}
\makeatother

\setcounter{secnumdepth}{5}
% Note: Toc changed for appendex
\setcounter{tocdepth}{1}
\definecolor{keywordcolor1}{rgb}{0,0,.4}
\definecolor{keywordcolor2}{rgb}{.90,0,0}

% See https://github.com/modelica-tools/listings-modelica/blob/master/listings-modelica.cfg
% Note: Changed comment color from green to [rgb]{0,0.4,0} - since the other variant was too distracting
% And added pure,impure,stream as keywords
% Remove cross as red
% Note: have basicstyle=\upshape\small\ttfamily in all languages,
% except the dialect [short]modelica - which is used for inline listings.
% Apart from that change that dialect should be identical to modelica

% The red keyword part is a bit unclear.
% The current logic is that "operators" are marked specially, but not normal functions.
%
% Note that Integer is both a predefined type and a function; we cannot treat them differently.
% Another solution would be to remove this completely.
\lstdefinelanguage{modelica}{% Language for use with the lstlisting environment.
    basicstyle=\upshape\small\ttfamily,        % Font size for displayed code listings.
    alsoletter={...},%
    %otherkeywords={-, =, +, [, ], (, ), \{, \}, :, *, !},%
    morekeywords=[1]{},% blue Keywords
    morekeywords=[2]{% blue + bold keywords
        annotation,assert,block,break,class,connector,constant,constrainedby,discrete,%
        each,encapsulated,else,elseif,elsewhen,end,exit,expandable,extends,external,final,flow,for,%
        function,if,in,inner,initial,input,import,loop,model,nondiscrete,operator,outer,%
        output,package,parameter,partial,record,redeclare,replaceable,return,%
        size,stream,terminate,then,type,when,while,algorithm,equation,%
        protected,public,and,false,not,or,true,pure,impure,transition,initialState},%
    % Note: initial is in both variants - depending on context, use first variant
    morekeywords=[3]{% red keywords - only special functions and classes: 3.7.2 #derivative-and-special-purpose-operators-with-function-syntax
    % 3.7.3 #event-related-operators-with-function-syntax (except initial) and 4.8 #predefined-types-and-classes
        connect,der,delay,cardinality,homotopy,semiLinear,inStream,actualStream,spatialDistribution,getInstanceName%
        terminal,noEvent,smooth,sample,pre,edge,change,reinit,%
        Real,Integer,Boolean,String},%
    comment=[l]{//}, % comment lines
    morecomment=[s]{/*}{*/}, % comment blocs
    morestring=[b]{'}, %
    morestring=[b]{"},
}[keywords,comments,strings]

\lstdefinelanguage[short]{modelica}[]{modelica}{% Language for use with the \lstinline command.
    basicstyle=\upshape\ttfamily,              % Font size for inline code snippets.
}

% Note: within only a keyword in grammar
\lstdefinelanguage{grammar}{%
    basicstyle=\upshape\small\ttfamily,        % size of fonts used for the code
    alsoletter={...},
    alsodigit={-},
    breaklines=true,
    breakatwhitespace=true,
    morekeywords=[2]{%
        letters,annotation,assert,block,break,class,connector,constant,constrainedby,discrete,%
        each,encapsulated,else,elseif,elsewhen,end,exit,expandable,extends,external,final,flow,for,%
        function,if,in,inner,initial,input,import,loop,model,nondiscrete,operator,outer,%
        output,package,parameter,partial,record,redeclare,replaceable,return,%
        size,stream,terminate,then,type,when,while,algorithm,equation,%
        protected,public,and,false,not,or,true,pure,impure,stream,transition,initialState,%
        within},
    morekeywords=[1]{|}
}[keywords,comments,strings]

\lstset{ %
  backgroundcolor=\color{white},   % choose the background color
  mathescape=true,
  breaklines=true,                 % automatic line breaking only at whitespace
  keepspaces,                      % don't remove space such as those after closing parenthesis
  captionpos=b,                    % sets the caption-position to bottom
  commentstyle=\color[rgb]{0,0.4,0}\sffamily,    % comment style
  keywordstyle=\color{blue}\ttfamily\bfseries,       % keyword style
  keywordstyle=[1]\color{keywordcolor1},
  keywordstyle=[2]\color{keywordcolor1}\bfseries,
  keywordstyle=[3]\color{keywordcolor2},
  stringstyle=\color{black},     % string literal style
  language=[short]modelica,     % Language that will be used by default, in particluar by plain \lstinline.
  showstringspaces=false,
  frame=lrtb, %
  xleftmargin=\fboxsep, %
  xrightmargin=-\fboxsep, %
  belowskip=0pt, %
}

% Duplicate this definition here to avoid issue - and name it "fortran77" to avoid problem in LatexML
\lstdefinelanguage{fortran77}{%
    basicstyle=\upshape\small\ttfamily,        % size of fonts used for the code
    morekeywords={%
        ASSIGN,BACKSPACE,CALL,CHARACTER,%
        CLOSE,COMMON,COMPLEX,CONTINUE,DATA,DIMENSION,DO,DOUBLE,%
        ELSE,ELSEIF,END,ENDIF,ENDDO,ENTRY,EQUIVALENCE,EXTERNAL,%
        FILE,FORMAT,FUNCTION,GO,TO,GOTO,IF,IMPLICIT,%
        INQUIRE,INTEGER,INTRINSIC,LOGICAL,%
        OPEN,PARAMETER,PAUSE,PRECISION,PRINT,PROGRAM,READ,REAL,%
        RETURN,REWIND,STOP,SUBROUTINE,THEN,%
        WRITE,SAVE},%
    morekeywords=[2]{%
        ACCESS,BLANK,BLOCK,DIRECT,EOF,ERR,EXIST,%
        FMT,FORM,FORMATTED,IOSTAT,NAMED,NEXTREC,NUMBER,OPENED,%
        REC,RECL,SEQUENTIAL,STATUS,TYPE,UNFORMATTED,UNIT},%
    morekeywords=[3]{%
        INT,DBLE,CMPLX,ICHAR,CHAR,AINT,ANINT,% left out real
        NINT,ABS,MOD,SIGN,DIM,DPROD,MAX,MIN,AIMAG,CONJG,SQRT,EXP,LOG,%
        LOG10,SIN,COS,TAN,ASIN,ACOS,ATAN,ATAN2,SINH,COSH,TANH,LGE,LLE,LLT,%
        LEN,INDEX},%
    morekeywords=[4]{AND,EQ,EQV,FALSE,GE,GT,OR,LE,LT,NE,NEQV,NOT,TRUE},%
    sensitive=f,%% not Fortran-77 standard, but allowed in Fortran-95 %%
    morecomment=[f]*,%
    morecomment=[f]C,%
    morecomment=[f]c,%
    morestring=[d]",%% not Fortran-77 standard, but allowed in Fortran-95 %%
    morestring=[d]'%
}[keywords,comments,strings]

\lstdefinelanguage{C}{%
    basicstyle=\upshape\small\ttfamily,        % size of fonts used for the code
    morekeywords={%
        auto,break,case,char,const,continue,default,do,double,%
        else,enum,extern,float,for,goto,if,int,long,register,return,%
        short,signed,sizeof,static,struct,switch,typedef,union,unsigned,%
        void,volatile,while},%
    sensitive,%
    morecomment=[s]{/*}{*/},%
    morecomment=[l]//,% nonstandard
    morestring=[b]",%
    morestring=[b]',%
    moredelim=*[directive]\#,%
    moredirectives={%
        define,elif,else,endif,error,if,ifdef,ifndef,line,%
        include,pragma,undef,warning}%
}[keywords,comments,strings,directives]

% Note: \lstinline!...! used for inline Modelica using [short]modelica

% Misc math
\makeatletter
\newcommand{\pdfrac}[3][\@nil]{%
  \def\tmp{#1}%
  \ifx\tmp\@nnil
    \frac{\mathrm{d}#2}{\mathrm{d}#3}%
  \else
    \frac{\mathrm{d}^{#1}#2}{\mathrm{d}#3^{#1}}%
  \fi%
}
\makeatother
\newcommand{\abs}[1]{\left\lvert #1{} \right\rvert}

% Text mode additions
\newcommand{\textgreatereq}{$\geq$}
\newcommand{\textlesseq}{$\leq$}

\usepackage[nameinlink,noabbrev]{cleveref}

% Change all autoref variants below chapter to use "section"
\def\subsectionautorefname{section}
\def\subsubsectionautorefname{section}
\def\paragraphautorefname{section}
\def\subparagraphautorefname{section}

% Make all description lists break the line after each key.
% henrikt-ma: Commenting out use of the enumitem package, see above.
%\setlist[description]{
%  style=nextline,
%  labelwidth=0pt,
%  leftmargin=15pt,
%  itemindent=\dimexpr-5pt-\labelsep\relax,
%}

\title{
\ifpdf
\begin{center}
\includegraphics[width=8cm]{Modelica_Language}
\end{center}
~\\[2\baselineskip]
\fi
Modelica\textsuperscript\textregistered - A Unified Object-Oriented Language for Systems
Modeling\\[2\baselineskip]
Language Specification\\[2\baselineskip]
Version 3.5-dev}
\date{\today}
\author{Modelica Association}
